\documentclass{amsart}
% For dealing with references we use the comment environment
\usepackage{verbatim}
\newenvironment{reference}{\comment}{\endcomment}
%\newenvironment{reference}{}{}
\newenvironment{slogan}{\comment}{\endcomment}
\newenvironment{history}{\comment}{\endcomment}

% For commutative diagrams we use Xy-pic
\usepackage[all]{xy}

% We use 2cell for 2-commutative diagrams.
\xyoption{2cell}
\UseAllTwocells

% We use multicol for the list of chapters between chapters
\usepackage{multicol}

% This is generally recommended for better output
\usepackage{lmodern}
\usepackage[T1]{fontenc}

% For cross-file-references

% Package for hypertext links:
\usepackage{hyperref}

% For any local file, say "hello.tex" you want to link to please

% Theorem environments.
%
\theoremstyle{plain}
\newtheorem{theorem}[subsection]{Theorem}
\newtheorem{proposition}[subsection]{Proposition}
\newtheorem{lemma}[subsection]{Lemma}

\theoremstyle{definition}
\newtheorem{definition}[subsection]{Definition}
\newtheorem{example}[subsection]{Example}
\newtheorem{exercise}[subsection]{Exercise}
\newtheorem{situation}[subsection]{Situation}

\theoremstyle{remark}
\newtheorem{remark}[subsection]{Remark}
\newtheorem{remarks}[subsection]{Remarks}

\numberwithin{equation}{subsection}

% Macros
%
\def\lim{\mathop{\mathrm{lim}}\nolimits}
\def\colim{\mathop{\mathrm{colim}}\nolimits}
\def\Spec{\mathop{\mathrm{Spec}}}
\def\Hom{\mathop{\mathrm{Hom}}\nolimits}
\def\Ext{\mathop{\mathrm{Ext}}\nolimits}
\def\SheafHom{\mathop{\mathcal{H}\!\mathit{om}}\nolimits}
\def\SheafExt{\mathop{\mathcal{E}\!\mathit{xt}}\nolimits}
\def\Sch{\mathit{Sch}}
\def\Mor{\mathop{\mathrm{Mor}}\nolimits}
\def\Ob{\mathop{\mathrm{Ob}}\nolimits}
\def\Sh{\mathop{\mathit{Sh}}\nolimits}
\def\NL{\mathop{N\!L}\nolimits}
\def\CH{\mathop{\mathrm{CH}}\nolimits}
\def\proetale{{pro\text{-}\acute{e}tale}}
\def\etale{{\acute{e}tale}}
\def\QCoh{\mathit{QCoh}}
\def\Ker{\mathop{\mathrm{Ker}}}
\def\Im{\mathop{\mathrm{Im}}}
\def\Coker{\mathop{\mathrm{Coker}}}
\def\Coim{\mathop{\mathrm{Coim}}}

% Boxtimes
%
\DeclareMathSymbol{\boxtimes}{\mathbin}{AMSa}{"02}

%
% Macros for moduli stacks/spaces
%
\def\QCohstack{\mathcal{QC}\!\mathit{oh}}
\def\Cohstack{\mathcal{C}\!\mathit{oh}}
\def\Spacesstack{\mathcal{S}\!\mathit{paces}}
\def\Quotfunctor{\mathrm{Quot}}
\def\Hilbfunctor{\mathrm{Hilb}}
\def\Curvesstack{\mathcal{C}\!\mathit{urves}}
\def\Polarizedstack{\mathcal{P}\!\mathit{olarized}}
\def\Complexesstack{\mathcal{C}\!\mathit{omplexes}}
% \Pic is the operator that assigns to X its picard group, usage \Pic(X)
% \Picardstack_{X/B} denotes the Picard stack of X over B
% \Picardfunctor_{X/B} denotes the Picard functor of X over B
\def\Pic{\mathop{\mathrm{Pic}}\nolimits}
\def\Picardstack{\mathcal{P}\!\mathit{ic}}
\def\Picardfunctor{\mathrm{Pic}}
\def\Deformationcategory{\mathcal{D}\!\mathit{ef}}

\setlength{\emergencystretch}{3em}
\begin{document}
\title[Stacks Project, Tag 0F3P]{565 --- Weightings on quasi-finite finitely presented affine morphisms descend through inverse limits}
\maketitle
\noindent Copyright (C) 2005--2025 Johan de Jong. Stacks Project authors. Source: \href{https://stacks.math.columbia.edu/tag/0F3P}{Tag 0F3P}.

\noindent Editorial difficulty note (not author text). Difficulty H2: Descending the constructible level sets does not by itself descend the weighting condition. The proof uses a universal finite open-and-closed test neighbourhood to encode that condition as local constancy, then descends openness of its finitely many level sets at a later stage..

\noindent Editorial provenance note (not author text). History: excerpt from revision \texttt{a0444\allowbreak 6e57e\allowbreak c1fbc\allowbreak 252a8\allowbreak 71afc\allowbreak ec775\allowbreak 2fb28\allowbreak 07b14}, more-morphisms.tex, lines 22890--22950. Reference presentation was adapted; original-author context and core support blocks were retained or added. The mathematical wording is unchanged. Licensed under GNU FDL 1.2 or later; no invariant sections or cover texts. See ../licenses/COPYING. Context blocks reproduce author definitions and supporting results; they are not additional counted problems.

\section{Weightings}
\label{section-weightings}

\noindent
The material in this section is taken from \href{https://stacks.math.columbia.edu/bibliography}{Bibliography: SGA4 (Exposee XVII, 6.2.4)}.

\medskip\noindent
Let $\pi : U \to V$ be a locally quasi-finite morphism of schemes with
finite fibres. Given a function $w : U \to \mathbf{Z}$ we define a function
$$
\textstyle{\int}_\pi w : V \longrightarrow \mathbf{Z},\quad
v \longmapsto
\sum\nolimits_{u \in U,\ \pi(u) = v} w(u) [\kappa(u) : \kappa(v)]_s
$$
Note that the field extensions are finite
(Morphisms, Lemma \href{https://stacks.math.columbia.edu/tag/01TG}{lemma residue field quasi finite (Tag 01TG)}),
$[\kappa' : \kappa]_s$ is the separable degree
(Fields, Definition \ref{fields-definition-insep-degree}),
and the sum is finite as the fibres of $\pi$ are assumed finite.
Another way to compute the value of $\int_\pi w$ at a point $v \in V$
is as follows.
Choose an algebraically closed field $k$ and a morphism
$\overline{v} : \Spec(k) \to V$ whose image is $v$. Then we have
$$
(\textstyle{\int}_\pi w)(v) =
\sum\nolimits_{\overline{u} \in U_{\overline{v}}} w(\overline{u})
$$
where of course $w(\overline{u})$ denotes the value of $w$ at the
image $u$ of the point $\overline{u}$ under the morphism
$U_{\overline{v}} \to U$.
Note that we may view $\overline{u} \in U_{\overline{v}}$ as morphisms
$\overline{u} : \Spec(k) \to U$ such that
$\pi \circ \overline{u} = \overline{v}$. Namely, since
$U \to V$ is locally quasi-finite with finite fibres,
the scheme $U_{\overline{v}}$ is the spectrum
of a finite dimension algebra over $k$ and all of whose prime ideals
are maximal ideals with residue field $k$. To see that the equality
holds, note that the number of morphisms $\overline{u}$ lying over
a given $u$ is equal to $[\kappa(u) : \kappa(v)]_s$ by
Fields, Lemma \href{https://stacks.math.columbia.edu/tag/09HJ}{lemma separable degree (Tag 09HJ)}.



\begin{definition}
\label{definition-weighting}
Let $f : X \to Y$ be a locally quasi-finite morphism. A
{\it weighting} or a {\it pond\'eration} of $f$ is a map
$w : X \to \mathbf{Z}$ such that for any diagram
$$
\xymatrix{
X \ar[d]_f & U \ar[l]^h \ar[d]^\pi \\
Y & V \ar[l]_g
}
$$
where $V \to Y$ is \'etale, $U \subset X_V$ is open, and $U \to V$ finite,
the function $\int_\pi (w \circ h)$ is locally constant.
\end{definition}

\begin{definition}
\label{fields-definition-insep-degree}
Let $E/F$ be an algebraic field extension. Let $E_{sep}$ be the subextension
found in Lemma \href{https://stacks.math.columbia.edu/tag/030K}{lemma separable first (Tag 030K)}.
\begin{enumerate}
\item The integer $[E_{sep} : F]$ is called the {\it separable
degree} of the extension. Notation $[E : F]_s$.
\item The integer $[E : E_{sep}]$ is called the {\it inseparable
degree}, or the {\it degree of inseparability} of the extension.
Notation $[E : F]_i$.
\end{enumerate}
\end{definition}

\begin{lemma}
\label{lemma-descend-weighting}
Let $Y = \lim Y_i$ be a directed limit of affine schemes. Let $0 \in I$
and let $f_0 : X_0 \to Y_0$ be a morphism of affine schemes which is
quasi-finite and of finite presentation. Let $f : X  \to Y$
and $f_i : X_i \to Y_i$ for $i \geq 0$ be the base changes of $f_0$.
If $w : X \to \mathbf{Z}$ is a weighting of $f$, then for sufficiently
large $i$ there exists a weighting $w_i : X_i \to \mathbf{Z}$
of $f_i$ whose pullback to $X$ is $w$.
\end{lemma}

\begin{proof}
By Lemma \href{https://stacks.math.columbia.edu/tag/0F3G}{lemma jumps w (Tag 0F3G)} the level sets of $w$ are constructible subsets
$E_k$ of $X$. This implies the function $w$ only takes a finite number
of values by Properties, Lemma
\href{https://stacks.math.columbia.edu/tag/0F2M}{lemma stratification locally finite constructible (Tag 0F2M)}.
Thus there exists an $i$ such that $E_k$ descends to a construcible
subset $E_{i, k}$ in $X_i$ for all $k$; moreover, we may assume
$X_i = \coprod E_{i, k}$. This follows as the topological space
of $X$ is the limit in the category of topological spaces
of the spectral spaces $X_i$ along a directed system with
spectral transition maps. See
Limits, Section \href{https://stacks.math.columbia.edu/tag/081A}{section descent (Tag 081A)}
and
Topology, Section \href{https://stacks.math.columbia.edu/tag/0A2U}{section limits spectral (Tag 0A2U)}.
We define $w_i : X_i \to \mathbf{Z}$ such that its level
sets are the constructible sets $E_{i, k}$.

\medskip\noindent
Choose $Y_{i, univ} \to Y_i$ and
$U_{i, univ} \subset Y_{i, univ} \times_{Y_i} X_i$
as in Lemma \href{https://stacks.math.columbia.edu/tag/0F3N}{lemma open and closed in quasi finite (Tag 0F3N)}.
By the universal property of the construction, in order to
show that $w_i$ is a weighting, it would suffice to show
that
$$
\tau_i = \textstyle{\int}_{U_{i, univ} \to Y_{i, univ}} w_i|_{U_{i, univ}}
$$
is locally constant on $Y_{i, univ}$. By Lemma \href{https://stacks.math.columbia.edu/tag/0F3H}{lemma jumps int w (Tag 0F3H)}
this function has constructible level sets but it may
not (yet) be locally constant. Set
$Y_{univ} = Y_{i, univ} \times_{Y_i} Y$
and let $U_{univ} \subset Y_{univ} \times_Y X$
be the inverse image of $U_{i, univ}$.
Then, since the pullback of $w$ to $Y_{univ} \times_Y X$
is a weighting for $Y_{univ} \times_Y X \to Y_{univ}$
(Lemma \href{https://stacks.math.columbia.edu/tag/0F3B}{lemma weighting base change (Tag 0F3B)})
we do have that
$$
\tau = \textstyle{\int}_{U_{univ} \to Y_{univ}} w_i|_{U_{univ}}
$$
is locally constant on $Y_{univ}$. Thus the level sets of
$\tau$ are open and closed. Finally, we have
$Y_{univ} = \lim_{i' \geq i} Y_{i', univ}$
and the level sets of $\tau$ are the inverse limits of the
level sets of $\tau_{i'}$ (similarly defined).
Hence the references above imply that for sufficiently
large $i'$ the level sets of $\tau_{i'}$ are open as well.
For such an index $i'$ we conclude that $w_{i'}$
is a weighting of $f_{i'}$ as desired.
\end{proof}
\end{document}
